\section{Uniform control of a whole small--prime face}
\label{sec:tpc44-robust-face}

Fix \(z\ge2\).  Every squarefree equality label has the unique split
\begin{equation}
 u=a b,
 \qquad
 \Pplus(a)\le z<\Pminus(b),
 \label{eq:tpc44-label-low-high}
\end{equation}
where \(a=1\) or \(b=1\) is allowed.  Let
\begin{equation}
 \mathcal A_z:=\{a:a\text{ occurs as the low part of an actual label}\},
 \qquad \nu_z:=|\mathcal A_z|.
 \label{eq:tpc44-low-projection-set}
\end{equation}
For a high--prime environment \(\eta=(\varepsilon_p)_{p>z}\), group
the coordinate coefficients by their low parts:
\begin{equation}
 B_{i,a}(\eta)
 :=\sum_{\substack{u\in\mathcal U_i\\u_{\le z}=a}}
 b_{i,u}\chi_{u_{>z}}(\eta),
 \qquad
 W_z(\eta):=\sum_i\sum_{a\in\mathcal A_z}|B_{i,a}(\eta)|^2.
 \label{eq:tpc44-B-W}
\end{equation}

\begin{lemma}[Projected diagonal identity]
\label{lem:tpc44-projected-diagonal}
One has
\begin{equation}
 \E_{\eta}W_z(\eta)=\cD_0.
 \label{eq:tpc44-projected-mean}
\end{equation}
\end{lemma}

\begin{proof}
Fix \(i\) and \(a\).  Two different full labels \(u,u'\in\mathcal U_i\)
with the same low part \(a\) have different high parts.  Indeed
\(u=ab=u'=ab'\) would imply \(u=u'\).  Their high Walsh characters are
therefore orthogonal.  Hence
\begin{equation}
 \E_\eta|B_{i,a}(\eta)|^2
 =\sum_{\substack{u\in\mathcal U_i\\u_{\le z}=a}}|b_{i,u}|^2.
 \label{eq:tpc44-B-diagonal}
\end{equation}
Summing over \(i\) and \(a\) gives \eqref{eq:tpc44-projected-mean}.
\end{proof}

\begin{theorem}[Robust face compression]
\label{thm:tpc44-robust-face}
For every \(L\ge1\),
\begin{equation}
 \Pp_\eta\!\left(
  \sup_{\sigma\in\{-1,1\}^{\{p\le z\}}}
  Z_0(\sigma,\eta)>L\nu_z\cD_0
 \right)
 \le L^{-1}.
 \label{eq:tpc44-robust-face-tail}
\end{equation}
Consequently there exists a high--prime environment \(\eta\) for which
\begin{equation}
 \sup_\sigma Z_0(\sigma,\eta)\le\nu_z\cD_0.
 \label{eq:tpc44-robust-face-exists}
\end{equation}
For each fixed small--prime assignment \(\sigma\), one also has
\begin{equation}
 \E_\eta Z_0(\sigma,\eta)\le\nu_z\cD_0.
 \label{eq:tpc44-fixed-face-mean}
\end{equation}
\end{theorem}

\begin{proof}
The grouping \eqref{eq:tpc44-B-W} gives
\begin{equation}
 S_i(\sigma,\eta)
 =\sum_{a\in\mathcal A_z}\chi_a(\sigma)B_{i,a}(\eta).
 \label{eq:tpc44-grouped-coordinate}
\end{equation}
Cauchy--Schwarz, uniformly in \(\sigma\), gives
\begin{equation}
 |S_i(\sigma,\eta)|^2
 \le \nu_z\sum_{a\in\mathcal A_z}|B_{i,a}(\eta)|^2.
 \label{eq:tpc44-coordinate-Cauchy}
\end{equation}
After summing over \(i\),
\begin{equation}
 \sup_\sigma Z_0(\sigma,\eta)\le\nu_zW_z(\eta).
 \label{eq:tpc44-uniform-W-bound}
\end{equation}
Markov's inequality and \cref{lem:tpc44-projected-diagonal} prove
\eqref{eq:tpc44-robust-face-tail}.  The finite average of \(W_z\)
equals \(\cD_0\), so some environment satisfies \(W_z\le\cD_0\); this gives
\eqref{eq:tpc44-robust-face-exists}.  Taking expectations in
\eqref{eq:tpc44-uniform-W-bound} proves
\eqref{eq:tpc44-fixed-face-mean}.
\end{proof}

No limiting or compactness argument is involved in the existential
statement.

The theorem includes a deterministic statement about the purely smooth
Walsh sector.  Averaging all high signs after fixing the small signs to
\(-1\) leaves exactly the kernels supported on primes at most \(z\).

\begin{corollary}[The smooth physical sector]
\label{cor:tpc44-smooth-sector}
One has
\begin{equation}
 0\le
 \cD_0+
 \sum_{\substack{\kappa>1\\\Pplus(\kappa)\le z}}
 \mu(\kappa)C(\kappa)
 \le\nu_z\cD_0.
 \label{eq:tpc44-smooth-sector}
\end{equation}
\end{corollary}

\begin{proof}
The middle expression is
\(
 \E_\eta Z_0(-\one_{\le z},\eta_{>z})
\).
It is nonnegative, and its upper bound is
\eqref{eq:tpc44-fixed-face-mean}.
\end{proof}

\begin{remark}[Why there is no union bound]
\label{rem:tpc44-no-union-bound}
The random variable \(W_z\) dominates all small--prime assignments at
once.  We therefore pay the actual projection count \(\nu_z\), not a
probability union over \(2^{\pi(z)}\) vertices.  Conditional moment
bounds for each fixed vertex would in any case require controlling
collisions after the low coordinates are removed; the grouping above
handles those collisions exactly.
\end{remark}
